\chapter{Results}
\label{ch:results}

All results in this chapter were produced by running the complete analysis
pipeline (using \code{./run\_circuit.sh --all}) against the fifteen valid and six
invalid benchmark circuits supplied with the assignment, on the codebase
version accompanying this report.

\section{Structural Validation on the Invalid Benchmark Set}
\label{sec:res-invalid}

All six intentionally invalid circuits were correctly rejected, each with
the specific error class the circuit was constructed to exercise
(\cref{tab:error-classes}). \Cref{tab:res-invalid} reproduces the exact
diagnostic message emitted for each.

\begin{table}[H]
\centering
\small
\begin{tabular}{@{}lp{0.62\textwidth}@{}}
\toprule
\textbf{Circuit} & \textbf{Reported diagnostic} \\
\midrule
\code{e01\_undefined\_cell} & \code{undefined cell 'MUX2\_X1'} \\
\code{e02\_undriven\_signal} & \code{signal 'X' is undriven} \\
\code{e03\_multiple\_drivers} & \code{net 'N1' has multiple drivers ('G1' and 'G2')} \\
\code{e04\_combinational\_loop} & \code{combinational loop involving gates: G1, G2, G3} \\
\code{e05\_wrong\_input\_count} & \code{cell 'NAND2\_X1' expects 2 inputs, got 1} \\
\code{e06\_undriven\_output} & \code{primary output 'OUT' is undriven} \\
\bottomrule
\end{tabular}
\caption{Validation results on the six invalid benchmark circuits.}
\label{tab:res-invalid}
\end{table}

\section{Nominal Timing on the Valid Benchmark Set}
\label{sec:res-nominal}

\Cref{tab:res-nominal} reports, for all fifteen valid circuits, the gate
count and the nominal (pre-optimization) circuit delay, WNS, and TNS. One of the fifteen circuits (\code{c01\_inverter\_chain}) is already timing-compliant
(nonnegative WNS); the remaining six exhibit a deliberate timing violation
by construction, ranging from a marginal \SI{-0.002}{ns} on
\code{c12\_monte\_carlo\_swit\newline ching} to a severe \SI{-0.147}{ns} (with a TNS
of \SI{-2.321}{ns}, reflecting many simultaneously violating output
transitions) on \code{c15\_high\_fanout\_sizing\_fabric}, the largest
circuit in the suite at 118 gates.

\newpage
\begin{longtable}{@{}lrrrr@{}}
\toprule
\textbf{Circuit} & \textbf{Gates} & \textbf{Delay (ns)} & \textbf{WNS (ns)} & \textbf{TNS (ns)} \\
\midrule
\endhead
c01\_inverter\_chain            & 3   & 0.1010 & \phantom{-}0.0090 & \phantom{-}0.0000 \\
c02\_basic\_nand\_path          & 3   & 0.1620 & -0.0070 & -0.0070 \\
c03\_fanout\_branch             & 4   & 0.2010 & -0.0060 & -0.0060 \\
c04\_reconvergent\_paths        & 5   & 0.2722 & -0.0072 & -0.0072 \\
c05\_multi\_output              & 5   & 0.2253 & -0.0053 & -0.0078 \\
c06\_deep\_critical\_path       & 10  & 0.4527 & -0.0227 & -0.0268 \\
c07\_parallel\_paths            & 6   & 0.2748 & -0.0078 & -0.0078 \\
c08\_non\_unate\_logic          & 4   & 0.2687 & -0.0087 & -0.0130 \\
c09\_high\_fanout               & 6   & 0.2090 & -0.0040 & -0.0073 \\
c10\_three\_input\_gates        & 5   & 0.2870 & -0.0120 & -0.0120 \\
c11\_gate\_sizing\_stress       & 6   & 0.3595 & -0.0295 & -0.0295 \\
c12\_monte\_carlo\_switching    & 7   & 0.3172 & -0.0022 & -0.0022 \\
c13\_large\_reconvergent\_network & 74  & 1.0103 & -0.0293 & -0.2133 \\
c14\_layered\_reconvergent\_fabric & 104 & 0.9893 & -0.0213 & -0.1596 \\
c15\_high\_fanout\_sizing\_fabric & 118 & 1.9674 & -0.1474 & -2.3211 \\
\bottomrule
\caption{Nominal (pre-optimization) timing on the fifteen valid benchmark
circuits. Nine of fifteen circuits are supplied already timing-compliant;
the remaining six exhibit a constructed violation.}
\label{tab:res-nominal}
\end{longtable}

\section{Gate-Sizing Heuristic Comparison}
\label{sec:res-heuristics}

\Cref{tab:res-heuristics} reports the final WNS reached by each of the
three automatically-invoked heuristics (\code{slack\_weighted\_capacitance},
the canonical logical-effort-guided method; \code{critical\newline ity\_effort\_gap},
the second logical-effort-informed heuristic; and \code{random\_greedy}, the
non-logical-effort baseline), together with a compliance flag indicating
whether that heuristic reached a fully timing-compliant final design within
the configured area and power budgets.

\begin{center}
    

\begin{longtable}{@{}lrrr@{}}
\toprule
\textbf{Circuit} & \textbf{WNS$_{\mathrm{swg}}$} & \textbf{WNS$_{\mathrm{ceg}}$} & \textbf{WNS$_{\mathrm{rg}}$} \\
\midrule
\endhead
c01\_inverter\_chain            & \phantom{-}0.0090 & \phantom{-}0.0090 & \phantom{-}0.0090 \\
c02\_basic\_nand\_path          & \phantom{-}0.0050 & \phantom{-}0.0050 & -0.0010  \\
c03\_fanout\_branch             & \phantom{-}0.0150 & \phantom{-}0.0150 & \phantom{-}0.0150 \\
c04\_reconvergent\_paths        & \phantom{-}0.0059 & \phantom{-}0.0059 & \phantom{-}0.0069 \\
c05\_multi\_output              & -0.0025 & -0.0025 & -0.0053 \\
c06\_deep\_critical\_path       & \phantom{-}0.0033 & \phantom{-}0.0033 & -0.0217 \\
c07\_parallel\_paths            & \phantom{-}0.0112 & \phantom{-}0.0112 & -0.0078 \\
c08\_non\_unate\_logic          & \phantom{-}0.0093 & \phantom{-}0.0093 & \phantom{-}0.0093 \\
c09\_high\_fanout               & \phantom{-}0.0430 & \phantom{-}0.0430 & \phantom{-}0.0430 \\
c10\_three\_input\_gates        & \phantom{-}0.0330 & \phantom{-}0.0330 & \phantom{-}0.0330 \\
c11\_gate\_sizing\_stress       & \phantom{-}0.0325 & \phantom{-}0.0325 & -0.0220 \\
c12\_monte\_carlo\_switching    & \phantom{-}0.0129 & \phantom{-}0.0129 & -0.0002 \\
c13\_large\_reconvergent\_network & \phantom{-}0.0037 & -0.0293 & -0.0293 \\
c14\_layered\_reconvergent\_fabric & \phantom{-}0.0182 & \phantom{-}0.0182 & \phantom{-}0.0030 \\
c15\_high\_fanout\_sizing\_fabric & \phantom{-}0.0392 & \phantom{-}0.0397 & \phantom{-}0.0397 \\
\midrule
\textbf{Repair rate} & \textbf{14/15} & \textbf{13/15} & \textbf{8/15} \\
\bottomrule
\caption{Final WNS (ns) and compliance status per heuristic. Repair rate
counts circuits reaching \(\mathrm{WNS} \geq 0\) within area/power
budgets, out of all fifteen valid circuits.}
\label{tab:res-heuristics}
\end{longtable}
\end{center}

Restricted to only the fourteen circuits that exhibited a violation to
begin with (excluding \code{c01}, already compliant), the repair rates are
$13/14$ ($\approx 92.9\%$) for slack-weighted capacitance, $12/14$
($\approx 85.7\%$) for criticality effort gap, and $7/14$ ($=50.0\%$) for
random greedy heuristic.

\subsection{Efficiency: STA Calls}
\label{sec:res-efficiency}

Summed across all fifteen circuits, the canonical logical-effort-guided
heuristic issued \textbf{119} total STA calls, versus \textbf{155} for
random greedy --- and the correct explanation for this is more specific
than ``one attempt, one STA call.'' \Cref{sec:opt-eval} already
establishes that candidate evaluation is a two-step gate: a cheap,
STA-free pre-filter (\code{replacement\_area\_and\_po\newline wer}) rejects a
candidate outright if it alone would exceed the area or power budget,
and only a candidate that survives this filter reaches the actual
\code{analyze\_timing} call. Cross-checking the recorded iteration logs
confirms this directly: on \code{c13\_large\_reconvergent\_network}, the
canonical heuristic's log contains 66 logged attempts, of which 52 are
rejected with reason \emph{``candidate violates area or power
constraints''} (free, no STA call) and only 13 require an actual timing
result (10 rejected for not improving WNS/TNS, 1 rejected once compliant
for not improving cost, and 2 accepted) --- plus one further STA call to
establish the initial reference state (\cref{sec:opt-cost}), for a
reported total of 14, exactly matching \code{summary.json}.

Because budget-infeasible candidates cost nothing, a heuristic's STA-call
total measures how many candidates it proposes that are \emph{plausible
enough to need a real timing check}, not merely how many candidates it
tries in total. On this specific circuit, random greedy needed only 3 STA
calls in total (1 initial, 1 accepted, 1 WNS/TNS rejection) against the
canonical heuristic's 14 --- the canonical heuristic's higher-scored
candidates more often clear the cheap budget pre-filter and therefore more
often reach the expensive check, while most of random greedy's random
draws are eliminated for free before ever costing a timing analysis. This
cautions against reading the suite-wide STA-call total as a simple
``search order wastes fewer attempts'' story: the totals do favor the
canonical heuristic in aggregate (119 versus 155) and as a per-circuit
mean ($\approx 7.9$ versus $\approx 10.3$ calls per circuit), but not
uniformly --- on five of the fifteen circuits, including the two largest
(\code{c13}: 14 versus 3; \code{c14}: 39 versus 19), the canonical
heuristic issues \emph{more} STA calls than random greedy, precisely
because more of its proposals are good enough to be worth the expensive
check in the first place. Aggregate STA-call count is therefore best read
as a rough proxy for total search cost, not as a clean, uniformly-favorable
efficiency metric on its own; \cref{tab:res-heuristics,tab:res-cost}'s
repair-rate and final-cost comparisons are the more direct evidence for
the canonical heuristic's advantage.

\subsection{Cost, Power, and Area}
\label{sec:res-cost}

\Cref{tab:res-cost} compares the final cost $J$ (\cref{eq:cost}) reached by
the canonical heuristic against random greedy. The canonical method
reaches a \emph{strictly} lower (better) final cost than random greedy on
ten of the fifteen circuits, and \emph{ties exactly} on the remaining
five: on \code{c01} because neither heuristic makes any accepted change at
all (the circuit is already compliant and no candidate reduces cost), and
on \code{c03}, \code{c08}, \code{c09}, and \code{c10} because both
heuristics independently accept the identical sequence of resizes (verified
directly against the recorded iteration logs, not merely inferred from
matching final cost). On \emph{no} circuit in the suite does random greedy
reach a lower final cost than the canonical heuristic. The narrowest
margin in the canonical method's favor is on \code{c15}
($J = 1.263$ versus $J = 1.272$), still a real, if small, improvement
rather than a tie.

\begin{longtable}{@{}lrrr@{}}
\toprule
\textbf{Circuit} & \textbf{$J_0$} & \textbf{$J_{\mathrm{swc}}$} & \textbf{$J_{\mathrm{rg}}$} \\
\midrule
\endhead
c01\_inverter\_chain            & 1.400 & 1.400 & 1.400 \\
c02\_basic\_nand\_path          & 1.616 & 1.511 & 1.578 \\
c03\_fanout\_branch             & 1.549 & 1.342 & 1.342 \\
c04\_reconvergent\_paths        & 1.532 & 1.434 & 1.481 \\
c05\_multi\_output              & 1.518 & 1.466 & 1.587 \\
c06\_deep\_critical\_path       & 1.650 & 1.377 & 1.659 \\
c07\_parallel\_paths            & 1.543 & 1.424 & 1.573 \\
c08\_non\_unate\_logic          & 1.561 & 1.376 & 1.376 \\
c09\_high\_fanout               & 1.496 & 1.207 & 1.207 \\
c10\_three\_input\_gates        & 1.609 & 1.368 & 1.368 \\
c11\_gate\_sizing\_stress       & 1.810 & 1.283 & 1.780 \\
c12\_monte\_carlo\_switching    & 1.434 & 1.412 & 1.444 \\
c13\_large\_reconvergent\_network & 1.545 & 1.376 & 1.551 \\
c14\_layered\_reconvergent\_fabric & 1.508 & 1.291 & 1.375 \\
c15\_high\_fanout\_sizing\_fabric & 1.775 & 1.263 & 1.272 \\
\bottomrule
\caption{Final cost function value $J$, initial versus
each heuristic's final design.}
\label{tab:res-cost}
\end{longtable}

\subsection{Case Study: \texttt{c02\_basic\_nand\_path}}
\label{sec:res-case-c02}

\Cref{fig:c02-pre,fig:c02-post} show the rendered pre- and
post-optimization circuit diagrams for the smallest circuit on which
slack-weighted capacitance and random greedy diverge in outcome, and the
full recorded iteration logs
(\code{optimization\_slack\_weighted\_capacitance.csv},
\code{optimization\_random\_greedy.csv}) make the reason for the divergence
precisely traceable rather than merely inferable from the final numbers.

Slack-weighted capacitance's very first candidate --- ranked highest by
its urgency$\times$mismatch score (\cref{sec:opt-swc}) --- is
\code{G3: AND2\_X1}$\rightarrow$\code{AND2\_X2}, which alone resolves the
violation ($\mathrm{WNS}: -0.0070 \rightarrow +0.0050$~ns) in one accepted
change. Three further candidates (\code{G1}, \code{G3}$\rightarrow$
\code{X4}, \code{G2}) are then proposed by the cost-reduction phase
(\cref{sec:opt-accept}) and all three are \emph{rejected} for violating the
configured area/power budget, at which point the search terminates
\code{no\_feasible\_change} --- a normal, compliant stop, not a failure.

Random greedy's uniformly random draw instead tries \code{G1:
NAND2\_X1}$\rightarrow$\code{X2} first (accepted; $\mathrm{WNS}:
-0.0070\rightarrow-0.0025$~ns, an improvement but not yet compliant), then
rejects \code{G3: AND2\_X1}$\rightarrow$\code{X2} for exceeding the area
budget (at this point, with \code{G1} already upsized, the combined area
would reach $10.2$, well past the $8.5$ limit), then accepts \code{G2:
INV\_X1}$\rightarrow$\code{X2} (WNS$\rightarrow-0.0010$~ns), and finally
exhausts its remaining candidates (\code{G1}$\rightarrow$\code{X4},
\code{G2}$\rightarrow$\code{X4}, a repeated \code{G3}$\rightarrow$\code{X2}
attempt) all rejected on the same area ceiling, terminating
\code{no\_feasible\_change} while \textbf{still in violation}.

Both searches plateau at the identical configured budget ceiling
(\code{maximum\_area} $= 8.5$), reached via two different combinations of
accepted resizes: slack-weighted capacitance reaches a total area of
$8.4$ after its single accepted \code{G3} resize alone, while random
greedy reaches the same $8.4$ via its two smaller accepted resizes
(\code{G1} then \code{G2}) combined. From that point, every further
candidate in \emph{either} run --- regardless of which gate it targets ---
is rejected against the same $8.5$ ceiling; the outcome is decided not by
which gate the budget happens to block, but by whether a compliant WNS was
already reached \emph{before} the ceiling was hit. Slack-weighted
capacitance's score-ranked draw tried the one resize that both fixes the
violation and fits the pristine budget on its very first attempt; random
greedy's random draw spent that same budget across two resizes that
together do not fix the violation, before any further candidate could be
tried. This is a concrete, single-circuit illustration of exactly the
aggregate effect quantified in \cref{sec:res-efficiency}: an informed
search order does not expand what a heuristic is allowed to do, only the
likelihood that a fix is found before the shared budget ceiling is
reached by changes that do not resolve the violation.

\begin{figure}[H]
    \centering
    \includegraphics[width=0.92\textwidth]{c02_pre.png}
    \caption{\texttt{c02\_basic\_nand\_path}, pre-optimization. Nodes are
    colored by slack; the critical path is highlighted.}
    \label{fig:c02-pre}
\end{figure}

\begin{figure}[H]
    \centering
    \includegraphics[width=0.92\textwidth]{c02_post.png}
    \caption{\texttt{c02\_basic\_nand\_path}, post-optimization (canonical
    slack-weighted-capacitance heuristic). Gate \texttt{G3} has been
    resized from \texttt{AND2\_X1} to \texttt{AND2\_X2}, resolving the
    timing violation.}
    \label{fig:c02-post}
\end{figure}

\Cref{fig:c02-convergence} makes the same story visible directly as a
convergence trace: slack-weighted capacitance (blue) reaches zero TNS and
positive WNS in a single accepted step and holds there, while random
greedy (green) needs two accepted steps to make partial progress and never
crosses zero. \Cref{fig:c02-outcomes} summarizes the same run as final
outcomes --- note criticality effort gap (orange) is visually
indistinguishable from slack-weighted capacitance on this particular
circuit, since both happen to find the same \texttt{G3} fix, consistent
with \cref{tab:res-heuristics}.

\begin{figure}[H]
    \centering
    \includegraphics[width=\textwidth]{c02_basic_nand_path_conv.png}
    \caption{\texttt{c02\_basic\_nand\_path}: normalized cost, WNS, TNS,
    circuit delay, area, and power versus accepted-change iteration, for
    all three sequential heuristics (\cref{sec:opt-heuristics}). The
    horizontal dashed line marks the compliance threshold on the WNS/TNS
    panels.}
    \label{fig:c02-convergence}
\end{figure}

\begin{figure}[H]
    \centering
    \includegraphics[width=\textwidth]{c02_basic_nand_path_outc.png}
    \caption{\texttt{c02\_basic\_nand\_path}: final runtime, total
    iterations (accepted and rejected attempts combined), and final WNS
    per heuristic.}
    \label{fig:c02-outcomes}
\end{figure}

\subsection{Case Study: Divergence Between the Two Logical-Effort Heuristics}
\label{sec:res-case-divergence}

\code{c13\_large\_reconvergent\_network} (74 gates) is the one circuit in
the suite on which slack-weighted capacitance and criticality effort gap
disagree, and the recorded iteration logs make the mechanism traceable
precisely rather than merely observable in the final numbers.

Criticality effort gap's search accepts exactly one change: at iteration 7,
\code{G65: OR2\_X1}\newline $\rightarrow$\code{OR2\_X2} is accepted under the
\emph{tied-WNS, TNS-improves} clause of the timing-repair acceptance rule
(\cref{sec:opt-accept}) --- WNS is left exactly unchanged at
$-0.0293$~ns, while TNS improves substantially, from $-0.2133$ to
$-0.1717$~ns. Eleven iterations later, at iteration 18, it proposes
\code{G7: NOR3\_X1}$\rightarrow$\code{NOR3\_X2} --- a resize that,
critically, \emph{is} available to criticality effort gap, since it is a
plain one-step-up candidate, not something requiring the logical-effort
bracket --- but this candidate is now \textbf{rejected} with reason
\emph{``candidate violates area or power constraints''}: the area consumed
by the earlier \code{G65} change has left no further budget for it. No
subsequent candidate improves WNS or TNS either, and the search terminates
\code{no\_feasible\_change} with the circuit still at its original
$-0.0293$~ns violation.

Slack-weighted capacitance's search proposes that exact same
\code{G7: NOR3\_X1}$\rightarrow$\code{NOR3\_X2} resize much earlier, at
iteration 13, while the circuit is still at its pristine initial area
budget (its own higher-scored earlier candidates were all rejected
outright rather than accepted, so no budget had yet been spent). At that
point it is well within budget and is accepted, cutting the violation's
magnitude by \textbf{88.6\%} ($-0.0293\rightarrow-0.0033$~ns); a second
accepted change at
iteration 35, \code{G6: NAND3\_X1}$\rightarrow$\code{NAND3\_X2}, then
closes the remaining gap entirely ($\mathrm{WNS}\rightarrow+0.0037$~ns,
$\mathrm{TNS}\rightarrow 0$).

This is a second, independent confirmation of the same effect identified
in \cref{sec:res-case-c02}: the outcome is decided not by what a heuristic
is theoretically permitted to consider --- both heuristics have \code{G7}
available --- but by whether an earlier accepted change quietly consumes
budget that a later, more consequential change will need. In this instance
the effect is sharper than in the c02 case study, because
criticality effort gap's early accepted change is not even a step toward
resolving the violation directly (it is WNS-neutral, accepted purely for a
TNS improvement); slack-weighted capacitance's higher-scored ranking
happens to defer any budget-consuming change until the one that actually
matters, avoiding the trap entirely on this circuit.

\Cref{fig:c13-convergence} shows this divergence directly: the orange
(criticality effort gap) trace visibly plateaus at a still-negative WNS
after its single early step, its cost curve flattening well above the blue
(slack-weighted capacitance) trace, which instead makes a late but decisive
descent to full compliance once its own higher-ranked candidates ahead of
\texttt{G7} are exhausted.

\begin{figure}[H]
    \centering
    \includegraphics[width=\textwidth]{c13_large_reconvergent_network_conv.png}
    \caption{\texttt{c13\_large\_reconvergent\_network}: convergence trace
    for all three sequential heuristics. Criticality effort gap's single
    early, WNS-neutral step (\cref{sec:res-case-divergence}) is visible as
    the orange trace's premature plateau above zero WNS.}
    \label{fig:c13-convergence}
\end{figure}

\section{Monte Carlo Statistical Analysis: \texttt{c12\_monte\_carlo\_switching}}
\label{sec:res-montecarlo}

\Cref{tab:res-mc} reports the paired Monte Carlo comparison
(\cref{ch:monte-carlo}) for the one benchmark circuit with statistical
analysis enabled, using $M = 3000$ samples, $\sigma_g = 0.03$,
$\sigma_l = 0.05$, and random seed $1405$.

\begin{table}[H]
\centering
\small
\begin{tabular}{@{}lrr@{}}
\toprule
\textbf{Metric} & \textbf{Pre-optimization} & \textbf{Post-optimization} \\
\midrule
Mean circuit delay (ns) & 0.3194 & 0.3041 \\
Std.\ dev.\ circuit delay (ns) & 0.0124 & 0.0116 \\
Mean WNS (ns) & -0.0044 & \phantom{-}0.0109 \\
Violation probability & 0.631 & 0.174 \\
Timing yield & 0.369 & 0.826 \\
Modal critical path (probability) & \parbox[t]{4.1cm}{\raggedright \texttt{D->G4->G5->G6->G7} (0.606)} & \parbox[t]{4.1cm}{\raggedright \texttt{A->G1->G2->G3->G7} (0.628)} \\
\bottomrule
\end{tabular}
\caption{Paired Monte Carlo statistical timing comparison, pre- versus
post-optimization, \texttt{c12\_monte\_carlo\_switching}, $M=3000$ samples.}
\label{tab:res-mc}
\end{table}

Optimization improves timing yield from \textbf{36.9\% to 82.6\%} --- a
$+45.7$ percentage-point improvement --- confirming that the
nominal-corner WNS improvement reported in \cref{sec:res-heuristics} for
this circuit translates into a substantial improvement in robustness to
manufacturing variation, not merely a nominal-case artifact. Notably, the
\emph{identity} of the modal critical path also changes: under process
variation, the pre-optimization design is most often limited by the path
through \code{G4--G7} (driven by primary input \code{D}), while the
post-optimization design is most often limited by an entirely different
path through \code{G1--G3, G7} (driven by primary input \code{A}) --- the
sizing changes did not merely widen the margin on the original critical
path, they shifted which path is critical at all, which is exactly the
kind of effect a single nominal-corner STA run cannot reveal.

\section{Synthetic Benchmark Suite Results}
\label{sec:res-benchmarking}
 
This section reports results from \code{complex\_repair\_suite}, a 60-case
synthetic benchmark generated and evaluated with the framework of
\cref{ch:benchmarking}. Unlike the fixed suite of \cref{sec:res-heuristics},
which compares three heuristics on fifteen circuits (eleven of them with only
a single-digit gate count), this suite is specifically configured to stress
the optimizer on large, structurally diverse circuits, and evaluates all
\textbf{four} heuristics --- including brute force, absent from the fixed
suite's reported comparison (\cref{sec:opt-bruteforce}) --- against the
independent counterfactual oracle of \cref{sec:bench-oracle}.
 
\subsection{Suite Composition and Generation Reliability}
\label{sec:res-bench-generation}

The suite comprises 60 cases: 40 procedurally generated (structured-random
topologies, 150--250 gates, 8--20 logic levels, targeting a
$0.35$ reconvergence probability and up to 6-way fan-out) and 20 seeded
from \code{c15\_high\_fanout\_sizing\_fabric} (\cref{sec:res-nominal}) under
independently sampled electrical conditions and planted mutations. Every
case's reference assignment was independently certified by an 8-restart,
beam-width-64 multi-start beam search (\code{best\_known\_multi\_start\_beam},
\cref{sec:bench-reference}) before any downsizing mutation was planted, and
each resulting initial (post-mutation) circuit was verified to have a
negative WNS strictly within the configured area/power budget.
 
Generation reliability was total for this run: of 60 requested cases, all
60 were generated successfully, with \textbf{zero} entries in
\code{generation\_failures.csv} --- every case succeeded on its first
generation attempt, with no discarded partial attempts of any kind.
 
\subsection{Wilson-Interval Repair Rate}
\label{sec:res-bench-repair-rate}
 
\Cref{tab:bench-repair-rate} reports the case-level repair rate for each
heuristic with a 95\% Wilson confidence interval (\cref{sec:bench-metrics}).
Brute force and random greedy are run once and 20 times per case
respectively (the latter to characterize its stochastic variance, per
\cref{sec:bench-parallel}); the two logical-effort-informed heuristics,
being deterministic given a fixed circuit and candidate order
(\cref{sec:opt-heuristics}), are run once per case.
 
\begin{table}[H]
\centering
\small
\begin{tabular}{@{}lrrrl@{}}
\toprule
\textbf{Heuristic} & \textbf{Runs} & \textbf{Successes} & \textbf{Repair rate} & \textbf{95\% Wilson CI} \\
\midrule
Brute force & 60 & 60 & 100.0\% & [93.98\%, 100.00\%] \\
Criticality effort gap & 60 & 60 & 100.0\% & [93.98\%, 100.00\%] \\
Slack-weighted capacitance & 60 & 59 & 98.3\% & [91.14\%, 99.71\%] \\
Random greedy & 1200 & 1134 & 94.5\% & [93.06\%, 95.65\%] \\
\bottomrule
\end{tabular}
\caption{Case-level repair rate with 95\% Wilson confidence intervals (60 cases; random greedy evaluated with 20
repetitions per case, 1200 total runs).}
\label{tab:bench-repair-rate}
\end{table}
 
\begin{figure}[H]
    \centering
    \includegraphics[width=\textwidth]{benchmarking/heuristic_overview.png}
    \caption{Benchmark heuristic overview: repair rate with 95\% bootstrap
    confidence intervals, per-run wall-clock runtime distribution, mean STA
    calls per run (log scale), and final normalized cost among repaired
    cases only. Note the four-order-of-magnitude STA-call gap between
    brute force ($\approx 1835$ calls/run) and criticality effort gap
    ($\approx 37$ calls/run).}
    \label{fig:bench-heuristic-overview}
\end{figure}
 
Three results are worth drawing out beyond the headline repair rates.
First, criticality effort gap matches brute force's 100\% repair rate on
this suite exactly, and does so at $\approx 37$ mean STA calls per run
versus brute force's $\approx 1835$ --- a two-order-of-magnitude reduction
in search cost for an identical repair outcome on this population. Second,
slack-weighted capacitance --- the heuristic found strictly more reliable
than criticality effort gap on the fixed suite (\cref{sec:res-heuristics},
where it alone repaired \code{c13\_large\_reconvergent\_network}) --- is
here the \emph{less} reliable of the two logical-effort-informed
heuristics, with a single failure among its 60 cases. This is a genuine,
population-dependent reversal, not a contradiction: the fixed suite's
one divergence case (\cref{sec:res-case-divergence}) is a single 74-gate
circuit, and this suite's 60 cases, averaging $150$--$250$ gates, are a
different and much larger sample; a report should not extrapolate either
suite's heuristic ranking to circuits outside its own tested population.
Third, random greedy's 1200-run aggregate repair rate (94.5\%) sits
noticeably below its \emph{best-run} behavior implied by
\cref{fig:bench-case-reliability} below --- most of its failures are
concentrated in specific cases, not spread uniformly across the suite.
 
The single slack-weighted-capacitance failure is traceable to a specific
case: \code{case\_0032} (198 gates, 18 logic levels, 7 primary inputs, 8
primary outputs, initial $\mathrm{WNS}=-0.0357$~ns, reference
$\mathrm{WNS}=+0.0051$~ns). Both brute force ($\mathrm{WNS}=+0.0028$~ns)
and criticality effort gap ($\mathrm{WNS}=+0.0051$~ns, matching the
reference exactly) repair this case; slack-weighted capacitance terminates
\code{no\_feasible\_change} at $\mathrm{WNS}=-0.0003$~ns --- a residual
violation of under half a picosecond, not a large or qualitative failure,
but a failure by the strict binary repair-success criterion nonetheless.
 
\subsection{Oracle-Measured Selection Accuracy and Regret}
\label{sec:res-bench-oracle}
 
\Cref{tab:bench-oracle} reports, for every accepted decision each
heuristic made, how it compares to the independent counterfactual oracle
of \cref{sec:bench-oracle}: whether the chosen gate was the oracle's
single best-scoring gate at that state (\emph{gate hit rate}), whether the
chosen (gate, cell) pair was the oracle's single best exact move
(\emph{exact-move hit rate}), the median and worst-case timing regret in
WNS relative to the best move the oracle could show was available, and how
often the heuristic's final repaired cost beat the independently-certified
reference assignment's own cost.
 
\begin{table}[H]
\centering
\small
\begin{tabular}{@{}lrrrrr@{}}
\toprule
\textbf{Heuristic} & \textbf{Gate hit} & \textbf{Exact-move} & \textbf{Median regret} & \textbf{Worst regret} & \textbf{Beats ref.} \\
 & \textbf{rate} & \textbf{hit rate} & \textbf{(WNS, ns)} & \textbf{(WNS, ns)} & \\
\midrule
BF & 7.9\% & 7.7\% & 0.0031 & 0.0365 & 100.0\% \\
CEG & 4.1\% & 2.8\% & 0.0058 & 0.0803 & 85.0\% \\
SWC & 1.6\% & 1.4\% & 0.0159 & 0.2030 & 100.0\% \\
RG & 2.3\% & 1.5\% & 0.0210 & 0.2009 & 68.9\% \\
\bottomrule
\end{tabular}
\caption{Oracle-measured decision quality (BF = brute force, CEG = criticality effort gap,
SWC = slack-weighted capacitance, RG = random greedy).
Regret is the WNS gap between the heuristic's chosen move and the oracle's
best available move at that decision point, evaluated over every state a
heuristic actually visited (\cref{sec:bench-oracle}); lower is better.}
\label{tab:bench-oracle}
\end{table}
 
\begin{figure}[H]
    \centering
    \includegraphics[width=\textwidth]{benchmarking/selection_quality.png}
    \caption{Left: gate-selection quality against the oracle's single best
    gate/exact move, and precision/recall of each heuristic's changed
    gates against the originally planted downsizing mutations. Right:
    median WNS regret against the exhaustive one-move oracle.}
    \label{fig:bench-selection-quality}
\end{figure}
 
The gate- and exact-move hit rates are low in absolute terms for
\emph{every} heuristic, including brute force ($7.9\%$/$7.7\%$) --- this
is expected, not a sign of a broken oracle: the oracle's ``best gate'' is
computed over \emph{every} legal same-family replacement on \emph{every}
gate in the circuit (\cref{sec:bench-oracle}), while even brute force only
proposes candidates on gates lying on a currently-reported critical path,
within a logical-effort-derived capacitance bracket
(\cref{sec:opt-bruteforce}); a circuit this large routinely has several
gates capable of producing an equally-good or negligibly-worse timing
improvement, so an exact match to the oracle's single top pick is a strict
criterion. The more informative signal is the \emph{median regret} column:
brute force's median regret (0.0031~ns) is roughly half of
criticality effort gap's (0.0058~ns), which is in turn roughly a third of
slack-weighted capacitance's (0.0159~ns) and random greedy's (0.0210~ns)
--- an ordering that does \emph{not} match the repair-rate ordering of
\cref{tab:bench-repair-rate} (where criticality effort gap tied brute force
and beat slack-weighted capacitance). Repair rate measures whether a
heuristic \emph{eventually} reaches compliance at all; median regret
measures how close to optimal its individual decisions are along the way.
The two metrics are related but distinct, and a heuristic can rank well on
one without ranking identically on the other.
 
The precision/recall of each heuristic's changed gates against the
\emph{originally planted} mutated gates (\cref{sec:bench-planted}) shows a
striking, and initially counter-intuitive, pattern: criticality effort gap
and random greedy --- the two \emph{weaker} heuristics by repair rate and
regret --- have substantially \emph{higher} planted-gate precision
($65.8\%$ and $56.2\%$) than brute force and slack-weighted capacitance
($4.9\%$ and $10.9\%$), whose repairs touch the originally-mutated gates
only rarely. This is consistent with, not contradictory to, the framework's
own design rationale (\cref{sec:bench-oracle}): a compliant, budget-respecting
repair is very rarely unique, and a search that draws from a wider
candidate bracket (brute force, slack-weighted capacitance) is more likely
to find an equally valid fix at a \emph{different} gate than the one
originally downsized, while a search restricted to smaller, one-step-up
moves (criticality effort gap, random greedy) is more likely to end up
reversing something closer to the original mutation simply because fewer
alternative routes are available to it. This is precisely why
\cref{ch:benchmarking} treats exact recovery of the planted mutation as a
secondary diagnostic rather than the primary correctness signal: by that
narrower measure alone, the two least effective heuristics on this suite
would appear the most accurate.
 
Finally, \emph{beats reference cost} shows a clean split along the same
line: brute force and slack-weighted capacitance --- both drawing from the
full logical-effort capacitance bracket, the same candidate source
available to the beam-search reference itself (\cref{sec:bench-reference})
--- match or beat the reference's own final cost on \emph{every}
successful repair (100.0\%), while criticality effort gap and random
greedy, restricted to one-step-up moves, fall short of it on $15\%$ and
$31\%$ of their successful repairs respectively.
 
\subsection{Convergence and Runtime Scaling}
\label{sec:res-bench-convergence}
 
\begin{figure}[H]
    \centering
    \includegraphics[width=\textwidth]{benchmarking/repair_convergence.png}
    \caption{Fraction of all 1380 runs with timing repaired, as a function
    of total attempted optimizer iterations (accepted and rejected
    attempts combined). Runs that never repair remain flat at their final
    value indefinitely.}
    \label{fig:bench-repair-convergence}
\end{figure}
 
\Cref{fig:bench-repair-convergence} shows the same repair-rate story as
\cref{tab:bench-repair-rate} unfolding over search effort rather than as a
single endpoint number: criticality effort gap and brute force both reach
their full $100\%$ repair rate within about 10 iterations, slack-weighted
capacitance climbs more gradually and plateaus just under $99\%$, and
random greedy's much shallower climb --- still visibly rising at 100
iterations --- reflects the combined effect of its unranked candidate
order (\cref{sec:opt-random-greedy}) and the cases it never resolves at
all within any of its repetitions.

\begin{figure}[H]
    \centering
    \includegraphics[width=\textwidth]{benchmarking/runtime_scalability.png}
    \caption{Median optimizer runtime per case versus gate count (log
    scale), with a fitted trend line per heuristic.}
    \label{fig:bench-runtime-scalability}
\end{figure}
 
Runtime separates the four heuristics by roughly an order of magnitude at
each tier (\cref{fig:bench-heuristic-overview}, top right; mean and
standard deviation in \cref{tab:bench-runtime}): brute force averages
$4.48\pm2.96$~s per run; slack-weighted capacitance, $0.281\pm0.063$~s;
random greedy, $0.114\pm0.099$~s; and criticality effort gap, the fastest,
$0.093\pm0.076$~s. \Cref{fig:bench-runtime-scalability} shows brute
force's runtime is roughly flat with circuit size over this suite's
150--250 gate range (its cost is dominated by the size of the candidate
bracket it evaluates per iteration, not raw gate count), while the three
sequential heuristics' runtimes visibly \emph{decrease} with gate count ---
counter-intuitive at first glance, but consistent with the suite's
generation process: larger generated circuits in this suite do not
necessarily have proportionally more or deeper critical-path violations to
repair, so the number of accepted-plus-rejected attempts needed does not
scale directly with gate count either.
 
\begin{table}[H]
\centering
\small
\begin{tabular}{@{}lrrrr@{}}
\toprule
\textbf{Heuristic} & \textbf{Mean runtime (s)} & \textbf{Std.\ dev.\ (s)} & \textbf{Mean STA calls} & \textbf{Mean final cost} \\
\midrule
BF & 4.479 & 2.962 & 1835.3 & 1.291 \\
SWC & 0.281 & 0.063 & 98.6 & 1.329 \\
RG & 0.114 & 0.099 & 51.7 & 1.369 \\
CEG & 0.093 & 0.076 & 37.0 & 1.356 \\
\bottomrule
\end{tabular}
\caption{Runtime, search cost, and mean final cost among repaired cases. BF = brute force, CEG = criticality effort gap,
SWC = slack-weighted capacitance, RG = random greedy.}
\label{tab:bench-runtime}
\end{table}
 
\subsection{Per-Case Reliability}
\label{sec:res-bench-case-reliability}
 
\begin{figure}[H]
    \centering
    \includegraphics[width=\textwidth]{benchmarking/case_reliability.png}
    \caption{Top: per-case repair probability for a representative subset
    of cases, ordered by source and initial WNS (brute force and
    criticality effort gap are uniformly green --- always successful; the
    single dark-red cell is slack-weighted capacitance's one failure,
    \texttt{case\_0032}). Bottom: random greedy's repair probability
    (across its 20 repetitions) versus initial violation severity, marker
    size indicating gate count, split by source type.}
    \label{fig:bench-case-reliability}
\end{figure}
 
\Cref{fig:bench-case-reliability}'s lower panel makes the population-level
pattern underlying \cref{sec:res-bench-strata}'s severity stratification
directly visible at the case level: random greedy's repair probability is
$100\%$ for the great majority of cases regardless of source, but drops
into the $70$--$90\%$ range for a specific subset of harder cases, with
the two clearest outliers (both at $40\%$) occurring at moderate, not
extreme, violation severity --- reinforcing that repair difficulty for a
one-step-up, unranked search depends on circuit structure as much as on
how large the initial violation is.
 
\subsection{Stratified Breakdown}
\label{sec:res-bench-strata}
 
\Cref{tab:bench-strata} breaks the repair rate down by three groupings
computed automatically by the evaluation framework
(\cref{sec:bench-metrics}): circuit depth, reconvergence level, and
initial violation severity. (This suite's cases fall into only two buckets
each for depth and reconvergence --- \emph{medium}/\emph{deep} and
\emph{medium}/\emph{high} respectively --- reflecting the generation
configuration's target ranges, \cref{sec:res-bench-generation}; no
\emph{shallow} or \emph{low-reconvergence} cases were generated in this
run.)
 
\begin{table}[H]
\centering
\small
\begin{tabular}{@{}llrrrrr@{}}
\toprule
\textbf{Grouping} & \textbf{Stratum} & \textbf{Cases} & \textbf{BF} & \textbf{CEG} & \textbf{SWC} & \textbf{RG} \\
\midrule
\multirow{2}{*}{Depth} & Medium & 10 & 100.0\% & 100.0\% & 100.0\% & 94.0\% \\
 & Deep & 50 & 100.0\% & 100.0\% & 98.0\% & 94.6\% \\
\multirow{2}{*}{Reconvergence} & Medium & 5 & 100.0\% & 100.0\% & 100.0\% & 100.0\% \\
 & High & 55 & 100.0\% & 100.0\% & 98.2\% & 94.0\% \\
\multirow{3}{*}{Violation severity} & Mild & 8 & 100.0\% & 100.0\% & 100.0\% & 100.0\% \\
 & Moderate & 34 & 100.0\% & 100.0\% & 97.1\% & 95.4\% \\
 & Severe & 18 & 100.0\% & 100.0\% & 100.0\% & 90.3\% \\
\bottomrule
\end{tabular}
\caption{Repair rate by circuit depth, reconvergence level, and initial
violation severity. BF = brute force, CEG = criticality effort gap,
SWC = slack-weighted capacitance, RG = random greedy.}
\label{tab:bench-strata}
\end{table}
 
Brute force and criticality effort gap are unaffected by any stratum in
this suite, repairing $100\%$ of cases regardless of depth, reconvergence,
or severity. Random greedy is the only heuristic whose repair rate degrades
meaningfully with reconvergence (from $100\%$ at medium reconvergence to
$94.0\%$ at high) and with violation severity (from $100\%$ on mild
violations down to $90.3\%$ on severe ones) --- exactly where a
logical-effort-free, unranked search is least likely to stumble onto the
gate a severe or highly-reconvergent violation actually needs. Slack-weighted
capacitance shows a smaller version of the same pattern on depth
($100\%\rightarrow98.0\%$) and reconvergence ($100\%\rightarrow98.2\%$),
consistent with \code{case\_0032} --- its one failure --- being both a
deep and a reconvergent circuit. Somewhat counter-intuitively,
slack-weighted capacitance's repair rate on \emph{severe} violations
($100\%$) is \emph{higher} than on \emph{moderate} ones ($97.1\%$): the
harder-looking severity axis is not the one that predicts its failure on
this suite, depth and reconvergence are. This is exactly the kind of
heuristic-specific, non-obvious failure pattern the fixed fifteen-circuit
suite is too small to reveal, and is the clearest evidence in this report
that a large, structurally-varied synthetic suite adds information the
fixed suite alone cannot provide.


\section{Complete Per-Circuit Figures}
\label{sec:res-full-figures}

\Cref{sec:res-case-c02,sec:res-case-divergence} showed the pre-/post-
optimization circuit diagrams and convergence/outcome plots for the two
circuits most useful for illustrating \emph{why} the heuristics diverge.
The identical set of four figures --- pre-optimization diagram,
post-optimization diagram, convergence trace, and outcomes summary --- is
generated for every one of the fifteen valid benchmark circuits;
Appendix \ref{app:full-figures} reproduces the complete set for systematic
reference.

\section{Summary}
\label{sec:res-summary}

Across the fixed benchmark suite, the logical-effort-guided canonical
heuristic (slack weighted capacitance) resolved timing violations on
$13/14$ originally-violating circuits, versus $7/14$ for the
logical-effort-free random-greedy baseline, while simultaneously issuing
fewer total STA calls (119 versus 155) and reaching a strictly lower final
cost on ten of fifteen circuits (tying on the remaining five, and never
losing on any). The second logical-effort-informed
heuristic, criticality effort gap, performs comparably to the canonical
method on all but one circuit, where its different ranking signal misses a
repair that slack-weighted capacitance finds. The one Monte
Carlo-instrumented circuit shows that the nominal-case timing improvement
from optimization corresponds to a substantial real improvement in
statistical timing yield, and that the identity of the critical path itself
can shift as a result of sizing changes. These fixed-suite results are
necessarily based on a small number of circuits; \cref{sec:res-benchmarking}
reports the corresponding, statistically stronger comparison on the
synthetic benchmark suite once completed.
